% Chapter 7: Conclusion
% Target: 2-3 pages (500-800 words)

\section{Summary of Contributions}

This thesis presented a comprehensive comparative study of machine learning and large language model approaches to quantitative trading. Our key contributions include:

\begin{enumerate}
    \item \textbf{Dual-Track Framework}: We designed and implemented a systematic framework enabling fair comparison between ML and LLM paradigms, with careful attention to reproducibility and avoiding look-ahead bias.

    \item \textbf{Empirical Findings}: Through extensive backtesting on CSI300 and QQQ markets, we demonstrated that LLM-based strategies achieve superior risk-adjusted returns (Sharpe ratio 0.66 vs 0.17 on CSI300), while ML models offer better cost-efficiency and lower latency.

    \item \textbf{Interpretability Analysis}: We provided dual interpretability through SHAP attribution for ML models and natural language reasoning for LLMs, addressing the ``black-box trust crisis'' in quantitative finance.

    \item \textbf{Open-Source Implementation}: We released a complete, reproducible implementation including trained models, backtest results, and evaluation scripts.
\end{enumerate}

\section{Key Insights}

Our findings reveal fundamental trade-offs between the two paradigms:

\begin{description}
    \item[Performance vs Cost] LLMs achieve better returns but at higher computational cost.
    \item[Fitting vs Reasoning] ML models fit historical patterns; LLMs reason about market dynamics.
    \item[Speed vs Interpretability] ML offers speed; LLMs offer explainability.
\end{description}

The choice between paradigms should be guided by specific application requirements:

\begin{itemize}
    \item For \textbf{low-frequency trading} where explainability matters, LLMs are preferred.
    \item For \textbf{high-frequency trading} requiring low latency, ML models remain essential.
    \item For \textbf{risk management} during market stress, a hybrid approach combining LLM veto with ML signals may be optimal.
\end{itemize}

\section{Final Remarks}

The emergence of large language models represents a paradigm shift in quantitative finance. While traditional machine learning approaches have dominated the field for decades, LLMs offer a new paradigm based on semantic reasoning rather than pattern fitting.

Our work provides the first systematic evidence that this paradigm shift can lead to improved trading performance. However, we also demonstrate that the two paradigms are not mutually exclusive---the future likely lies in hybrid systems that combine the strengths of both.

As the field continues to evolve, we anticipate further improvements in both paradigms: more efficient LLMs with lower latency, and more sophisticated ML models with better generalization. The Dual-Track framework we propose provides a foundation for evaluating these advances in a fair and rigorous manner.

\vspace{2em}

\noindent\textit{``The best way to predict the future is to create it.''} --- Peter Drucker